\documentclass[../../../include/open-logic-section]{subfiles}

\begin{document}

\olfileid{sth}{spine}{foundation}

\olsection{Landhesan}

Kita mung \emph{meh} rampung---durung \emph{temenan}
rampung---amarga ing $\ZFminus$ ora ana kang
njamin manawa \emph{saben} himpunan kalebu ing
sawijining $V_\alpha$, yaiku saben himpunan
kabentuk ing sawijining tataran.

Saiki, miturut salah siji pangerten matematis kang
cukup prasaja, kita ora \emph{perlu nggatekake}
apa ana himpunan ing sanjabane hierarki. (Yen ana,
kita bisa wae nglirwakake.) Nanging, \emph{konsep}
kita babagan himpunan dilandhesi gagasan yen saben
himpunan kabentuk ing sawijining tataran
(delengen \stageshier{} ing \olref[z][story]{sec}).
Mula kita kudu nyingkirake kamungkinan anane
himpunan ing sanjabane hierarki. Kanggo kuwi,
kita kudu nambah aksioma anyar kang njamin yen
saben himpunan ana ing salah siji tataran hierarki.

Amarga $V_\alpha$ iku tataran-tataran kita,
kita bisa nganggep pratelan iki minangka aksioma:

\begin{defish}
\emph{Regularitas.} $\forall A \exists \alpha\, A \subseteq V_\alpha$
\end{defish}

Iki pendekatan kang bisa ditampa. Nanging,
amarga alasan kang bakal diterangake ing bagean
sabanjure, kita milih aksioma alternatif:

\begin{axiom}[Landhesan]
$(\forall A \neq \emptyset)(\exists B \in A)A \cap B = \emptyset$.
\end{axiom}

Kanthi pambuktian kang rada dawa, kita bisa
nuduhake (ing $\ZFminus$) yen Landhesan
ngakibatake Regularitas:
\begin{defn}
Kanggo saben himpunan $A$, tetepna:
\begin{align*}
	\text{cl}_0(A) &= A,\\
	\text{cl}_{n+1}(A) &= \bigcup \text{cl}_n(A),\\
	\text{trcl}(A) &= \bigcup_{n < \omega} \text{cl}_{n}(A).
\end{align*}
$\text{trcl}(A)$ diarani \emph{panutupan
transitif} saka $A$.
\end{defn}
\noindent
Jeneng ``panutupan transitif'' iku trep:
\begin{prop}\ollabel{subsetoftrcl}
$A \subseteq \trcl{A}$ and $\trcl{A}$ is a transitive set. 
\end{prop}

\begin{proof}
Cetha yen $A = \text{cl}_0(A) \subseteq
\text{trcl}(A)$. Yen $x \in b \in \trcl{A}$,
mula $b \in \text{cl}_n(A)$ kanggo sawijining
$n$, saengga $x \in \text{cl}_{n+1}(A)
\subseteq \text{trcl}(A)$.
\end{proof}

\begin{lem}\ollabel{lem:TransitiveWellFounded}
Yen $A$ himpunan transitif, ana sawijining
$\alpha$ kang nyukupi $A \subseteq V_\alpha$.
\end{lem}

\begin{proof}
Elinga definisi ``$\supstrict(X)$'' saka
\olref[ordinals][opps]{defsupstrict}. Tetepna
loro himpunan iki:
\begin{align*}
	D  &= \Setabs{x \in A}{\forall \delta\ x \nsubseteq V_\delta}\\
	\alpha &= \supstrict\Setabs{\delta}{(\exists x \in A)
	(x \subseteq V_\delta \land (\forall \gamma \in \delta)x \nsubseteq V_\gamma)}
\end{align*}
Upamane $D = \emptyset$. Dadi yen $x \in A$,
ana sawijining $\delta$ kang nyukupi
$x \subseteq V_\delta$ lan, amarga ordinal
terurut becik, uga nyukupi
$(\forall \gamma \in \delta)x \nsubseteq V_\gamma$;
mula $\delta \in \alpha$ lan
$x \in V_\alpha$ miturut
\olref[spine][Valphabasic]{Valphabasicprops}.
Kanthi mangkono, $A \subseteq V_{\alpha}$,
kaya kang dikarepake.

Dadi cukup mbuktekake $D = \emptyset$.
Kanggo reductio, upamane kosokbaline. Miturut
Landhesan, ana sawijining $B \in D \subseteq A$
kang nyukupi $D \cap B = \emptyset$.
Yen $x \in B$, mula $x \in A$ amarga $A$
transitif; lan amarga $x \notin D$, mula
$\exists \delta\ x \subseteq V_\delta$.
Saiki tetepna
\[
	\beta = \supstrict\Setabs{\delta}{(\exists x \in B)(x \subseteq V_\delta \land (\forall \gamma < \delta)x \nsubseteq V_\gamma)}.
\]
Kaya sadurunge, $B \subseteq V_\beta$,
bertentangan karo $B \in D$.
\end{proof}

\begin{thm}\ollabel{zfentailsregularity}
Regularitas lumaku.
\end{thm}

\begin{proof}
Tetepna $A$; $A \subseteq \trcl{A}$ miturut
\olref{subsetoftrcl}, lan himpunan iki transitif.
Mula ana sawijining $\alpha$ kang nyukupi
$A \subseteq \trcl{A} \subseteq V_\alpha$
miturut \olref{lem:TransitiveWellFounded}.
\end{proof}

Asil-asil iki nuduhake yen $\ZFminus$ mbuktekake
implikasi $\emph{Landhesan}\Rightarrow\emph{Regularitas}$.
Ing \olref[spine][rank]{zfminusregularityfoundation},
kita bakal mbuktekake yen $\ZFminus$ uga mbuktekake
$\emph{Regularitas}\Rightarrow\emph{Landhesan}$.
Dadi, Landhesan lan Regularitas iku
\emph{ekivalen} (relatif marang $\ZFminus$).
Mula, yen $\ZFminus$ wis ditampa, kita bisa
menehi dhasar kanggo Landhesan liwat
ekivalensine karo Regularitas. Dene Regularitas
bisa langsung diwenehi dhasar saka
\stageshier{}.

\end{document}
